\section{第6章：静电场中的导体与电介质（6.5练习题）}
\sourceproblem{选择1}{99}{105}
\problemtext 判断静电平衡带电导体的电势性质。
\workings 导体内部 $\bm E=0$，连通导体为等势体，表面同势，选D。
\crosscheck 曲率影响局部电荷密度，不意味着电势不同。
\sourceproblem{选择2}{99}{105}
\problemtext 正带电无限平面面密度 $\sigma$，旁放平行中性导体板，求近远两面感应面密度。
\workings 设近面 $\sigma_1$、远面 $\sigma_2$，中性给 $\sigma_1+\sigma_2=0$。板内场为零给 $\sigma+\sigma_1-\sigma_2=0$，得 $-\sigma/2,+\sigma/2$，选B。
\crosscheck 默认板孤立中性且无额外外场；若接地或另有净电荷，条件改变。
\sourceproblem{选择3}{99}{105}
\problemtext 半径R的接地导体球外距心2R有点电荷q，求感应净电荷。
\workings 镜像法等效球内像电荷 $q'=-qR/(2R)=-q/2$，位置距心 $R/2$，外部场满足球面零势；由包围球的高斯面知感应净电荷为 $-q/2$，选D。
\crosscheck 接地球可交换电荷，与孤立中性球不同。
\sourceproblem{选择4}{99}{105}
\problemtext 相距很远的半径R及r金属球用细导线连接，求面密度比。
\workings 忽略相互影响及导线电荷，等势给 $Q_R/R=Q_r/r$，故 $\sigma_R/\sigma_r=r/R$，选D。
\crosscheck 电荷量与半径成正比，面密度与半径成反比。
\sourceproblem{选择5}{99}{105}
\problemtext 孤立导体球半径R、电势 $U_0$，求球外轴距r处场强大小（选项续PDF106）。
\workings $V(r)=RU_0/r$，故 $E=R|U_0|/r^2$，按正电势约定选C。
\crosscheck 场强大小须非负，若 $U_0<0$ 方向向内。
\sourceproblem{选择6}{100}{106}
\problemtext 同轴导体内半径a、接地外筒半径b，内线密度正 $\lambda$，求间隙r处场与势。
\workings $E=\lambda/(2\pi\varepsilon_0r)$，$V=\int_r^b E(s)ds=\lambda\ln(b/r)/(2\pi\varepsilon_0)$，选D。
\crosscheck $V(b)=0$，负梯度恢复场。
\sourceproblem{选择7}{100}{106}
\problemtext 均匀带电无限厚导体板两侧等距离h处，比较电势差。
\workings 无额外外场的对称模型中两侧场强大小相同，且板为等势体，因此距板两侧相等距离处电势相等，选A。
\crosscheck 若外场破坏对称性，则不能仅凭距离相同得出等势。
\sourceproblem{选择8}{100}{106}
\problemtext 判断介质中高斯定理的正确推论。
\workings $\oint\bm D\cdot d\bm S=Q_{\rm free,enc}$，净通量只由包围的自由电荷代数和确定，选C。
\crosscheck 面上D为零仅能推出净自由电荷为零，不能排除等量异号自由电荷。
\sourceproblem{选择9}{100}{106}
\problemtext 导体球外充满相对介电常量 $\varepsilon_r$ 的均匀介质，测表面附近场强E，求自由面电荷密度。
\workings 取外法向，$D_n=\sigma_f$，若题设正电荷，则 $\sigma_f=\varepsilon_0\varepsilon_r E$，选B。
\crosscheck 一般有向式 $\sigma_f=\varepsilon_0\varepsilon_r\bm E\cdot\bm n$；大小用绝对值。
\sourceproblem{选择10}{100}{106}
\problemtext 空气平行板电容器中平行插介质板，比较介质场E和空气场 $E_0$。
\workings 忽略边缘效应且界面无自由电荷，法向D连续，$E=E_0/\varepsilon_r$；普通介质 $\varepsilon_r>1$，方向相同，选C。
\crosscheck 这里是沿场方向分层，并非并排介质的同电压情形。
\sourceproblem{选择11}{100}{106}
\problemtext 点电荷旁非对称放置介质，对以电荷为中心的球面判断高斯定理适用性。
\workings 高斯定理仍成立，但非对称极化破坏球对称，单凭积分定理不能推出球面逐点场强，选B。
\crosscheck $\bm E$ 通量涉及全部电荷，$\bm D$ 通量涉及自由电荷；二者均不因对称性不足而失效。
\sourceproblem{选择12}{101}{107}
\problemtext 恒压平行板电容器填满均匀介质，比较E和D。
\workings 固定板距，$E=U/d=E_0$，$D=\varepsilon_0\varepsilon_rE=\varepsilon_rD_0$，选B。
\crosscheck 电源补充自由电荷；恒压与断电情形不同。
\sourceproblem{选择13}{101}{107}
\problemtext 面积S的两平行金属板分别带 $q_1,q_2$，间距d，求板间电势差。
\workings 忽略边缘效应，取A指向B为正，间隙场 $E=(q_1-q_2)/(2\varepsilon_0 S)$，故 $U_{AB}=(q_1-q_2)d/(2\varepsilon_0 S)$。若问大小，选C。
\crosscheck 原选项C带绝对值，仅是势差大小；有向 $U_{AB}$ 需保留符号。
\sourceproblem{选择14}{101}{107}
\problemtext 等外半径孤立空心和实心金属球的电容比较。
\workings 无内置独立电荷且外部介质相同，$C=4\pi\varepsilon_0R$，两者相等，选C。
\crosscheck 电容由外部几何及介质决定，不由金属体积决定。
\sourceproblem{选择15}{101}{107}
\problemtext 电容器充电断电后增大板距，比较电压、场强、储能。
\workings Q固定，$C=\varepsilon_0S/d$ 减小，$E=Q/(\varepsilon_0S)$ 不变，$U=Ed$ 及 $W=Q^2/(2C)$ 均增大，选C。
\crosscheck 外力拉开极板所做的功使储能增加；忽略边缘效应。
\sourceproblem{选择16}{101}{107}
\problemtext 固定空间电荷分布的体密度整体变为两倍，问电场能量倍数。
\workings 在固定线性介质、无不随之缩放的额外外场条件下，场整体加倍，$W=\int\varepsilon E^2dV/2$ 变为4倍，选C。
\crosscheck 若存在固定外场或非线性介质，不能直接使用此比例。
\sourceproblem{选择17}{101}{107}
\problemtext 空气平行板电容器断电后填满普通均匀介质，比较E、C、U、W。
\workings Q固定，$C'=\varepsilon_rC$，$E',U',W'$ 分别是原来的 $1/\varepsilon_r$；E减、C增、U减、W减，选B。
\crosscheck 假定 $\varepsilon_r>1$ 且几何不变。
\sourceproblem{选择18}{101}{107}
\problemtext 电场各点场强加倍，问电场能量倍数。
\workings 对固定线性介质及相同积分区域，$w=\varepsilon E^2/2$，总能量变为4倍，选B。
\crosscheck 若能量原本为零，则倍数概念退化；非线性介质不一定平方关系。
\sourceproblem{选择19}{101}{107}
\problemtext 无限大导体板两外表面面密度均为 $\sigma$，求外场大小。
\workings 每面贡献 $|\sigma|/(2\varepsilon_0)$，板外同向叠加，$E=|\sigma|/\varepsilon_0$，正密度约定选C。
\crosscheck 板内两场抵消；厚度不影响理想平面结果。
\sourceproblem{选择20}{101}{107}
\problemtext 同总电荷Q、同半径R的均匀带电实心球和球面，比较静电能（选项续PDF108）。
\workings 两者外部场完全相同，外区能量 $Q^2/(8\pi\varepsilon_0R)$。实心球内部 $E=Qr/(4\pi\varepsilon_0R^3)$，额外能量 $Q^2/(40\pi\varepsilon_0R)$，因此实心球总能量是球面的6/5，非零Q时较大，选B。
\crosscheck 球面内部场为零，外部能量相等而非一大一小。

\sourceproblem{填空1}{102}{108}
\problemtext 同心导体壳内壳带 $q$、外壳总电荷 $-2q$，求外壳内外表面电荷。
\workings 外壳金属内场为零，高斯面给内表面 $-q$；守恒给外表面 $-q$。
\crosscheck 两表面相加为 $-2q$，系统总电荷为 $-q$。
\sourceproblem{填空2}{102}{108}
\problemtext 负点电荷靠近孤立中性导体，比较导体内部电场及电势。
\workings 静电平衡内部场始终为零，不变；以无穷远为零势，导体电势降低。可用孤立导体的单位电势辅助解 $u$（导体表面1、无穷远0）及互易定理得导体势 $V=q_{\rm ext}u(\bm r_q)/C<0$。
\crosscheck 导体未接地且原先无其他场；接地时电势固定为零。
\sourceproblem{填空3}{102}{108}
\problemtext 导体外表面局部面电荷密度 $\sigma(x,y,z)$，求紧邻表面外侧电场。
\workings 真空中法向边界条件 $\bm E_{\rm out}=\sigma\bm n/\varepsilon_0$，故大小 $|\sigma|/\varepsilon_0$，正密度向外，负密度向内。
\crosscheck 切向场为零；公式不等于单独无限薄面的一半场。
\sourceproblem{填空4}{102}{108}
\problemtext 两远离的孤立球以细导线相连，问分配电荷与球半径关系。
\workings 等势给 $Q_A/R_A=Q_B/R_B$，电荷与半径成正比。
\crosscheck 忽略互相感应及导线电容。
\sourceproblem{填空5}{102}{108}
\problemtext 正点电荷q位于中性金属球外，球心至电荷矢径 $\bm r$，求感应净电荷及其在球心产生的场。
\workings 孤立中性球感应净电荷为0。外点电荷在球心产生 $-k_eq\bm r/r^3$，球内总场为零，因此感应场 $\bm E_{\rm ind}=k_eq\bm r/r^3$。
\crosscheck 感应净电荷零不意味着感应场零。
\sourceproblem{填空6}{102}{108}
\problemtext 真空中均匀带电导体球壳半径R、总电荷Q，求壳内电势。
\workings 以无穷远为零势，外区电势 $k_eQ/r$，由连续性内部恒为 $k_eQ/R$。
\crosscheck 空腔无独立电荷；内部场零。
\sourceproblem{填空7}{102}{108}
\problemtext 正负电荷中心原本重合的分子所属电介质类型及外场作用结果。
\workings 无极分子电介质；外场使中心相对位移，形成感生电偶极矩。
\crosscheck 与有极分子的取向极化区别。
\sourceproblem{填空8}{102}{108}
\problemtext 同轴圆筒间介质相对介电常量 $\varepsilon_r$，内外单位长度电荷 $\pm\lambda$，求间隙D和E大小。
\workings 圆柱高斯面给 $D=|\lambda|/(2\pi r)$，$E=|\lambda|/(2\pi\varepsilon_0\varepsilon_r r)$。
\crosscheck 适用于 $R_1<r<R_2$ 的长筒近似。
\sourceproblem{填空9}{102}{108}
\problemtext 真空中无限长均匀带正电线单位长度的电场通量。
\workings 取长度1 m的同轴闭合柱面，高斯定理给单位长度通量 $\lambda/\varepsilon_0$。
\crosscheck 电场线条数是绘图约定，严格物理量为通量每单位长度。
\sourceproblem{填空10}{102}{108}
\problemtext 内外半径 $R_1,R_2$ 的均匀带电介质球层总电荷Q，图中P位于介质内轴距r，求场强。
\workings 若Q为均匀分布的自由体电荷，$R_1<r<R_2$ 时 $Q_{\rm enc}=Q(r^3-R_1^3)/(R_2^3-R_1^3)$。故
\[
E_P=\frac{|Q|(r^3-R_1^3)}{4\pi\varepsilon_0\varepsilon_r r^2(R_2^3-R_1^3)}.
\]
\crosscheck 内界面处场为零；若题意Q含束缚电荷，则需区分总电荷与自由电荷，不能直接以同一公式处理。
\sourceproblem{填空11}{102}{108}
\problemtext 空气平板电容器恒压，原自由面电荷密度 $\pm\sigma$，填满介质后求E（续PDF109）。
\workings 电压及板距不变，$E=U/d=\sigma/\varepsilon_0$。
\crosscheck 新自由面密度为 $\varepsilon_r\sigma$，非原值。
\sourceproblem{填空12}{103}{109}
\problemtext 平板电容器充满均匀介质，自由面密度大小 $\sigma$，求D与E。
\workings $D=\sigma,\ E=\sigma/(\varepsilon_0\varepsilon_r)$。
\crosscheck 忽略边缘效应。
\sourceproblem{填空13}{103}{109}
\problemtext 半径R带电q金属球壳外充满无限均匀介质，内部真空，求球壳电势。
\workings 外部 $D_r=q/(4\pi r^2)$，积分得 $U=q/(4\pi\varepsilon_0\varepsilon_r R)$。
\crosscheck 取无穷远零势，内腔无电荷，外介质决定电容。
\sourceproblem{填空14}{103}{109}
\problemtext 同一对点电荷真空相距7 cm与煤油相距5 cm时力相等，求相对介电常量。
\workings $1/d_1^2=1/(\varepsilon_r d_2^2)$，故 $\varepsilon_r=49/25=1.96$。
\crosscheck 无限均匀线性介质模型；距离比需平方。
\sourceproblem{填空15}{103}{109}
\problemtext 板距d、初电压U的断电电容器插入厚 $d/3$ 的不接触极板金属片，求电压。
\workings 自由电荷固定，空气隙场不变，金属内E为零，总空气隙 $2d/3$，故 $U'=2U/3$。
\crosscheck 金属片不短接两板；忽略边缘效应，位置不影响总空气隙。
\sourceproblem{填空16}{103}{109}
\problemtext 同一对点电荷在真空距离 $r_1$ 与介质中距离 $r_2$ 时相互作用力相等，求相对介电常量。
\workings 库仑定律给 $\varepsilon_r=(r_1/r_2)^2$。
\crosscheck 假定电荷量不变且介质无限均匀。
\sourceproblem{填空17}{103}{109}
\problemtext 孤立导体球半径0.1 m、电势300 V，求距心0.3 m处电势。
\workings $V(r)=V(R)R/r=100$ V。
\crosscheck 无穷远零势；观察点位于球外。
\sourceproblem{填空18}{103}{109}
\problemtext 金属球壳半径R带电q，内腔填介质、外部真空，求壳电势。
\workings 无内置自由电荷时空腔场为零，外部场不受内介质影响，$U=q/(4\pi\varepsilon_0R)$。
\crosscheck 与填空13的外部填介质不同。
\sourceproblem{填空19}{103}{109}
\problemtext 孤立导体电荷q、电势U，写电容定义及意义。
\workings $C=q/U$，表征在单位电势下容纳电荷的能力。线性介质中C由几何和介质决定。
\crosscheck 取无穷远零势，$q=U=0$ 时由线性关系的比例系数定义，不能直接作0/0。
\sourceproblem{填空20}{103}{109}
\problemtext 平板电容器充电断电后增大板距，判断场强与电容。
\workings $E=Q/(\varepsilon_0 S)$ 不变，$C=\varepsilon_0S/d$ 减小。
\crosscheck 固定自由电荷、面积及介质，忽略边缘效应。
\sourceproblem{填空21}{103}{109}
\problemtext 雷雨云高500 m、面积 $10^7$ 平方米，均匀场 $10^4$ V/m，求全场能等效为从该高度落下的物体质量。
\workings $W=\varepsilon_0E^2Sh/2\approx2.2125\times10^6$ J；令 $mgh=W$，得 $m=\varepsilon_0E^2 S/(2g)\approx451.5$ kg（$g=9.8$）。
\crosscheck 这是均匀平行板储能估算，非实际闪电可释放能量；高度约去。
\sourceproblem{填空22}{103}{109}
\problemtext 图示电容器经电阻接恒压电源，插入充满介质后求电容及储能倍数。
\workings 最终稳态电阻无压降，电压不变，$C'=\varepsilon_r C_0,\ W'=C'U^2/2=\varepsilon_r W_0$。
\crosscheck 暂态电压可变化；题目问插入后的平衡状态。
\sourceproblem{填空23}{103}{109}
\problemtext 两空气电容器并联恒压电源，缓慢向其中一个插入普通介质，判断总电荷及储能。
\workings $C_{\rm eq}=C_1+C_2$ 增加，$Q_{\rm tot}=C_{\rm eq}U$ 与 $W=C_{\rm eq}U^2/2$ 均增大。
\crosscheck 恒压下另一电容器电荷不变，电源提供新增自由电荷。
\sourceproblem{填空24}{103}{109}
\problemtext 相对介电常量4、原题能量密度 $2\times10^6\,\mathrm{J/cm^3}$，求场强。
\workings 原页单位为每立方厘米，换算 $w=2\times10^{12}$ J每立方米。线性介质 $w=\varepsilon_0\varepsilon_r E^2/2$，故 $E=\sqrt{2w/(4\varepsilon_0)}\approx3.36\times10^{11}$ V/m。
\crosscheck 该数值远超普通介质的击穿场，原数据仅能作为形式计算；若单位原意为每立方米，结果为 $3.36\times10^8$ V/m，不能无依据替换原单位。
\sourceproblem{填空25}{104}{110}
\problemtext 平板电容器断电后填满均匀介质，问场强和储能倍数。
\workings Q不变，$E'=E/\varepsilon_r,\ W'=Q^2/(2\varepsilon_r C)=W/\varepsilon_r$。
\crosscheck 两个倍数相同，电容增大而储能降低。
\sourceproblem{计算1}{104}{110}
\problemtext 中间A板带 $3.0\times10^{-7}$ C，两侧接地B、C距A分别4.0、2.0 mm，面积均200平方厘米，求B、C感应电荷。
\workings 两间隙等效并联电容 $C_B=\varepsilon_0S/d_B,C_C=\varepsilon_0S/d_C$，$V_A=Q/(C_B+C_C)$。故 $Q_B=-Q C_B/(C_B+C_C)=-1.0\times10^{-7}$ C，$Q_C=-2.0\times10^{-7}$ C。
\crosscheck 感应总电荷 $-Q$；较近接地板电荷绝对值较大，忽略边缘外场。
\sourceproblem{计算2}{104}{110}
\problemtext 内导体球半径 $R_1$、电势 $V_0$，外同心薄壳半径 $R_2$、总电荷Q，求场势分布。
\workings 令内球电荷q，由 $V_0=k_e(q/R_1+Q/R_2)$ 得 $q=R_1(V_0/k_e-Q/R_2)$。外壳内外表面电荷为 $-q,Q+q$。取无穷远零势，
\[
V(r)=\begin{cases}V_0,&r\le R_1,\\ k_e(q/r+Q/R_2),&R_1<r<R_2,\\ k_e(q+Q)/r,&r\ge R_2.\end{cases}
\]
$E_r$ 依次为0、$k_eq/r^2$、$k_e(q+Q)/r^2$（界面取两侧极限）。
\crosscheck 势连续，导体金属内场零，负梯度恢复场。
\sourceproblem{计算3}{104}{110}
\problemtext 内球半径 $R_1$、电荷 $Q_1$，外壳内外半径 $R_2,R_3$、总电荷 $Q_2$，求电荷、势和场。
\workings 外壳内表面 $-Q_1$，外表面 $Q_1+Q_2$。令 $V_B=k_e(Q_1+Q_2)/R_3$，$V_A=V_B+k_eQ_1(1/R_1-1/R_2)$。径向场
\[
E_r=\begin{cases}0,&r<R_1,\\k_eQ_1/r^2,&R_1<r<R_2,\\0,&R_2<r<R_3,\\k_e(Q_1+Q_2)/r^2,&r>R_3.\end{cases}
\]
间隙势为 $V_B+k_eQ_1(1/r-1/R_2)$，两金属区分别为 $V_A,V_B$，外部势为 $k_e(Q_1+Q_2)/r$。
\crosscheck 内外表面电荷相加为 $Q_2$；所有界面电势连续。
\sourceproblem{计算4}{104}{110}
\problemtext 半径a、b的同心薄球壳，内壳带 $Q_1$，求令内壳电势零的外壳电荷及空间势。
\workings $0=k_e(Q_1/a+Q_2/b)$，故 $Q_2=-bQ_1/a$。取无穷远零势，
\[
V(r)=\begin{cases}0,&r\le a,\\k_eQ_1(1/r-1/a),&a<r<b,\\k_eQ_1(1-b/a)/r,&r\ge b.\end{cases}
\]
\crosscheck 在 $r=a,b$ 两侧连续，外部净电荷为 $Q_1(1-b/a)$，不必为零。

\sourceproblem{计算5}{105}{111}
\problemtext 内球半径 $R_1$ 带正电q，同心厚壳内外半径 $R_2,R_3$、原带正电Q，将壳内表面接地，求平衡电荷和内球电势。
\workings 整个连通壳均接地，内球q保持不变。壳内高斯面给内表面 $-q$；外部无独立电荷，零势边界唯一解为零，外表面电荷0。$V_1=k_eq(1/R_1-1/R_2)>0$。
\crosscheck 接地后外壳总电荷为 $-q$，初始Q不守恒，差额经地线交换。
\sourceproblem{计算6}{105}{111}
\problemtext 半径R接地球外距心r有点电荷q，求总感应电荷。
\workings 像电荷 $q'=-qR/r$ 位于源电荷方向距心 $R^2/r$ 处。球面任一点至像电荷距离为至真实电荷距离的 $R/r$ 倍，两势抵消。唯一性给外部解；包围球而不包围真实点电荷的高斯面给 $Q_{\rm ind}=q'$。
\crosscheck 接地球不等同于孤立中性球；$r\to\infty$ 时感应净电荷趋零。
\sourceproblem{计算7}{105}{111}
\problemtext 半径 $R_1<R_2$ 的同心薄金属球壳分别带 $Q_1,Q_2$，求场和球心势。
\workings 由球对称性，三个区域径向场分别为0、$k_eQ_1/r^2$、$k_e(Q_1+Q_2)/r^2$；球心势 $V(0)=k_e(Q_1/R_1+Q_2/R_2)$。
\crosscheck 壳面取场的两侧极限；电势连续且取无穷远零势。
\sourceproblem{计算8}{105}{111}
\problemtext 两绝缘同心薄球壳半径 $r_1<r_2$，起初中性，内壳加q；随后外壳接地再断开，最后内壳接地，求各阶段电荷与势。
\workings 第一阶段外壳净电荷0，内外表面为 $-q,+q$，$V_2=k_eq/r_2$。第二阶段外壳接地，外表面0，净电荷 $-q$，$V_2=0$。断开后外壳净电荷固定 $-q$。第三阶段内壳接地，设其电荷为q'，$0=k_e(q'/r_1-q/r_2)$，故 $q'=qr_1/r_2$；外壳电势变为 $k_eq(r_1/r_2-1)/r_2<0$，即相对第二阶段的变化。
\crosscheck 第三阶段外壳已绝缘，不可继续施加零势；内壳电荷减少 $q(1-r_1/r_2)$。
\sourceproblem{计算9}{106}{112}
\problemtext 半径0.5及1.0 m的两球相连，各自有内半径1.2 m接地同心壳；空气击穿场 $3\times10^6$ V/m，求首先击穿处和总电荷。
\workings 共同电势V，各球电容 $C_i=4\pi\varepsilon_0R_iR/(R-R_i)$。球表面场 $E_i=VR/[R_i(R-R_i)]$，两系数分别为 $24/7,6$（每米），故B球表面先达到阈值，$V=5.0\times10^5$ V。
\[
Q=(C_1+C_2)V=4\pi\varepsilon_0(6/7+6)(5\times10^5)\approx3.81\times10^{-4}\,\mathrm C.
\]
\crosscheck 各间隙场按 $1/r^2$ 减小，壳内表面场不超过内球表面；导线影响按原题忽略。
\sourceproblem{计算10}{106}{112}
\problemtext 均匀正带电环形薄片内外半径 $R_1,R_2$、面密度 $\sigma$，求轴上势及质子从无穷远穿越所需初速。
\workings 薄环积分
\[
V(x)=\frac{\sigma}{2\varepsilon_0}(\sqrt{x^2+R_2^2}-\sqrt{x^2+R_1^2}).
\]
势在中心最大 $V(0)=\sigma(R_2-R_1)/(2\varepsilon_0)$。能量守恒给临界速率 $v_c=\sqrt{e\sigma(R_2-R_1)/(m_p\varepsilon_0)}$。
\crosscheck 严格有限时穿过中心须 $v_0>v_c$，等号为渐近临界轨道。金属环均匀面密度按题设固定模型，实际孤立导体边缘电荷一般不均匀；非相对论速率式要求动能远小于静能。
\sourceproblem{计算11}{106}{112}
\problemtext 两接地无限板位于 $x=\pm d$，之间均匀正离子体密度 $\rho=nq$，求场势。
\workings 平面对称给 $E_x(0)=0$，$\partial_xE_x=\rho/\varepsilon_0$，故 $E_x=\rho x/\varepsilon_0$。积分并用 $V(\pm d)=0$，得 $V(x)=\rho(d^2-x^2)/(2\varepsilon_0)$。
\crosscheck 电势最大在中面；两板内侧自由面密度均为 $-\rho d$，单位面积总电荷守恒。
\sourceproblem{计算12}{107}{113}
\problemtext 孤立导体球半径R，逐步从无穷远搬运电荷，求已带q时增加dq的外力功及充至Q总功。
\workings 准静态 $dW=V(q)dq=k_eq\,dq/R$；积分 $W=k_eQ^2/(2R)=Q^2/(8\pi\varepsilon_0R)$。
\crosscheck 不可用最终电势乘全部电荷，否则多一倍；与场能积分一致。
\sourceproblem{计算13}{107}{113}
\problemtext 相距很远的半径1.0、2.0 cm球各带 $10^{-8}$ C，连接后求电荷及电势。
\workings 等势及总电荷守恒给 $Q_1=2\times10^{-8}/3$ C、$Q_2=4\times10^{-8}/3$ C。$V=k_e(2\times10^{-8})/(0.01+0.02)=6000$ V。
\crosscheck 忽略相互影响与导线电容；大球分得两倍电荷。
\sourceproblem{计算14}{107}{113}
\problemtext 带净电荷Q的厚球壳内外半径a,b，空腔中距心r有点电荷q，求壳面电荷、内表面在球心的势及球心总势。
\workings 内表面总电荷 $-q$，外表面 $Q+q$。内表面各元距心均为a，故 $V_{\rm in}(O)=-k_eq/a$。无外场时外球面均匀，$V_{\rm out}(O)=k_e(Q+q)/b$。对 $0<r<a$，
\[
V(O)=k_e[q/r-q/a+(Q+q)/b].
\]
\crosscheck 内表面虽一般不均匀，球心等距积分仍成立；若q位于中心，该点自身势奇异。
\sourceproblem{计算15}{107}{113}
\problemtext 同轴圆柱电容器半径 $R_1,R_2$，介质 $\varepsilon=\varepsilon_0\varepsilon_r$，自由线密度 $\pm\lambda$，求E、D、P、极化面密度与电容。
\workings $D_r=\lambda/(2\pi r)$，$E_r=D_r/\varepsilon$，$P_r=(1-1/\varepsilon_r)D_r$。介质内外边界外法向分别为负、正径向，故 $\sigma_{b,1}=-(1-1/\varepsilon_r)\lambda/(2\pi R_1)$，$\sigma_{b,2}=(1-1/\varepsilon_r)\lambda/(2\pi R_2)$。每单位长度电容 $C'=2\pi\varepsilon/\ln(R_2/R_1)$，长度L时 $C=LC'$（忽略端部）。
\crosscheck 介质内 $\nabla\cdot\bm P=0$，两束缚面电荷每米代数和为零。
\sourceproblem{计算16}{108}{114}
\problemtext 同轴电缆半径 $R_1,R_2$，介质分界半径R，内外两层相对介电常量分别 $\varepsilon_{r1},\varepsilon_{r2}$，求单位长度电容。
\workings 法向D在无自由面电荷的分界处连续，$D_r=\lambda/(2\pi r)$。两层径向势差相加，故
\[
C'=\frac{2\pi\varepsilon_0}{\ln(R/R_1)/\varepsilon_{r1}+\ln(R_2/R)/\varepsilon_{r2}}.
\]
\crosscheck 相同介质恢复均匀同轴公式；两层为串联关系而非并联。
\sourceproblem{计算17}{108}{114}
\problemtext 无限均匀介质介电常数 $\varepsilon$ 中半径R的孤立球充至Q，证明静电能公式。
\workings $V(q)=q/(4\pi\varepsilon R)$，准静态外功 $\int_0^Q V(q)dq=Q^2/(8\pi\varepsilon R)$。另由 $\int_R^\infty\varepsilon E^2(4\pi r^2)dr/2$ 得同值。
\crosscheck 线性无损介质、固定几何且可逆充电；$\varepsilon$ 为绝对介电常数。
\sourceproblem{应用1}{108}{114}
\problemtext 同轴真空二极管阴阳极半径0.50、4.50 mm，长度25 mm，阳极高300 V，求电子由静止到达阳极的动能、速率及阴极处电力。
\workings $K=eU=300$ eV，约 $4.81\times10^{-17}$ J；非相对论 $v=\sqrt{2eU/m_e}\approx1.03\times10^7$ m/s。长筒无空间电荷近似 $V(r)=V_1+U\ln(r/R_1)/\ln(R_2/R_1)$，$E_r=-U/[r\ln(R_2/R_1)]$，阴极电子所受力向外，$F=eU/[R_1\ln9]\approx4.37\times10^{-14}$ N。
\crosscheck 能量结果不依赖场分布；力的计算忽略端部效应和电子空间电荷。
\sourceproblem{应用2}{108}{114}
\problemtext 同轴电容器内外半径2、5 cm，均匀介质、恒压32 V，求轴距3.5 cm处场强及与外筒势差。
\workings 若内筒电势较高，$E_r=32/[0.035\ln(5/2)]\approx998$ V/m，$V_A-V_2=32\ln(5/3.5)/\ln(5/2)\approx12.46$ V。极性反转两有向量变号。
\crosscheck 恒压且均匀介质时场与相对介电常量无关；势差积分范围为A至外筒。
\sourceproblem{应用3}{109}{115}
\problemtext 二氧化钛晶片面积1.0平方厘米、厚0.10 mm，两侧紧贴极板，加12 V，求电容、自由和极化电荷、场强。
\workings 原题未给晶片介电常量，设沿法线有效相对介电常量 $\varepsilon_r$（线性模型）。$S/d=1$ m，$C=8.854\times10^{-12}\varepsilon_r$ F，$Q=1.0625\times10^{-10}\varepsilon_r$ C，$E=1.2\times10^5$ V/m，$|\sigma_f|=\varepsilon_0\varepsilon_r E\approx1.0625\times10^{-6}\varepsilon_r$ C每平方米，$|\sigma_b|=\varepsilon_0(\varepsilon_r-1)E$。
\crosscheck 二氧化钛的晶型和晶向等会影响介电响应，不擅自指定常数；束缚电荷在邻正自由电荷一侧为负。
\sourceproblem{应用4}{109}{115}
\problemtext 8、8、4微法电容器串联接12 V，求4微法上的电荷；断开电源后拆下再同极性并联，求电压。
\workings $C_s=(1/8+1/8+1/4)^{-1}=2$ 微法，故各电荷量24微库。并联前各正极电荷合计72微库，$C_p=20$ 微法，故 $U_f=3.6$ V。
\crosscheck 再连接时必须与电源隔离；能量初为144微焦、末为129.6微焦，差额耗散或辐射。
\sourceproblem{应用5}{109}{115}
\problemtext 板距 $d_0$、面积S的电容式测厚传感器，材料厚d、相对介电常量 $\varepsilon_r$，求C与d关系及金属情况。
\workings 忽略边缘且试样覆盖有效面积，空气与介质串联，
\[
C=\frac{\varepsilon_0S}{d_0-d+d/\varepsilon_r},\qquad
d=\frac{d_0-\varepsilon_0S/C}{1-1/\varepsilon_r}.
\]
绝缘于极板的金属试样内部场零，故 $C=\varepsilon_0S/(d_0-d)$。
\crosscheck $0\le d<d_0$，$\varepsilon_r=1$ 时无厚度灵敏度；试样须不短接两极，介电常量需标定。
\sourceproblem{应用6}{110}{116}
\problemtext 空气平板电容器正方形边长l、板距1 mm、100 V时电荷量 $10^{-4}$ C，求边长。
\workings $Q=\varepsilon_0 l^2 U/d$，故 $l=\sqrt{Qd/(\varepsilon_0U)}\approx10.63$ m。
\crosscheck 所需面积约113平方米，板距远小于边长，平板近似合理；实际大面积器件另有绝缘及结构限制。
\sourceproblem{应用7}{110}{116}
\problemtext 空气板距1.5 cm、电压40 kV、空气击穿场3 MV/m；插入厚0.30 cm、相对介电常量7.0、击穿场10 MV/m的玻璃，判断击穿。
\workings 原场 $40\,000/0.015=2.667$ MV/m，小于阈值。若维持40 kV，法向D连续，$E_a=U/(0.012+0.003/7)\approx3.218$ MV/m，$E_g=E_a/7\approx0.460$ MV/m，故空气先击穿而玻璃未达到其自身阈值。
\crosscheck 若插片前断电，固定自由电荷使空气场保持2.667 MV/m，玻璃场约0.381 MV/m，均不击穿；原题未明说是否保持电压，分别注明。
\sourceproblem{应用8}{110}{116}
\problemtext 电容键盘金属片面积50.0平方毫米，初距0.600 mm，检测电容增量0.250 pF，求按下距离。
\workings $C_0=\varepsilon_0S/d_0\approx0.73785$ pF，$C_f=C_0+0.250$ pF，$d_f=\varepsilon_0S/C_f$，故 $\Delta d=d_0-d_f\approx0.152$ mm。
\crosscheck 电容与间距成反比而非线性；求得位移小于初间距，忽略杂散电容变化与边缘效应。
